\section{Introduction}
\label{sec:introduction}

A visitor enters a virtual exhibition, points to a product display, and asks an
embodied agent, ``What is this?'' The environment might contain multiple agents
(such as a receptionist, a product specialist, or a career advisor), each with
unique responsibilities and capabilities. Generating a fluent answer does not
guarantee that the interaction was correct: the system must also track the
targeted object, identify the responsible agent, verify that the required
capability is explicitly declared for that agent, execute the intended action,
and ensure the resulting response remains anchored to the same object.

This execution chain easily breaks because its components are spread across
different architectural layers: a 3D engine (e.g., Unity) manages scene objects,
colliders, and user state; configuration files specify agent roles;
orchestration logic routes requests; and backend services call models or tools,
frequently logging nothing beyond request success. Consequently, individual
components may appear valid locally while the overall interaction fails---for
instance, the answer might refer to a different object, originate from an
unauthorized agent, or stem from an undeclared operation. Although our prototype
uses Unity, this cross-layer mismatch is representative of similar multi-tiered
architectures.

In intelligent user interfaces, successful network transport or model invocation
does not equal interaction validity. Before text generation starts, the system
must answer four deterministic questions: which object was targeted, which
agent owns responsibility for it, which capability is assigned to that agent
and request, and which specific backend action implements this capability.
Following generation, the system must verify that the output remains aligned
with the selected object and responsible agent.

We present an \emph{executable interaction contract}, \systemname{}
(Functional Meta-Language-defined Structure), extending the MLDS
(Meta-Language-defined Structure) artifact \cite{fischer2025machinereadable}.
This contract maintains alignment between target, provider, capability, action,
and evidence under a fixed model version, while rejecting and localizing
incomplete, outdated, ambiguous, unreachable, or conflicting execution chains.

Three key requirements, grounded in human--AI interaction guidelines
\cite{amershi2019guidelines,yang2020reexamining,weisz2024design}, drive this design:

\begin{description}
\item[Stable spatial reference.] The focused object remains the active
target throughout request construction, and the output is presented only when
returned identifiers match that target.

\item[Visible unresolved states.] Conditions like \emph{no target} or
\emph{ambiguous target} explicitly prevent deictic queries instead of
defaulting to a fluent yet ungrounded response.

\item[Capability-aware, inspectable routing.] Agent selection depends on
modeled roles and capabilities; the chosen provider, handoff step, capability,
and action are preserved as structured evidence to inform UI cues or user
explanations.
\end{description}

While these mechanisms establish a technical control boundary, their utility in
intelligent interfaces depends on whether users can comprehend the interaction
flow. The contract activates only after a target object is chosen; it neither
infers unstated intent nor guarantees the factual accuracy of generated text.
However, it guarantees that the chosen object is reliably linked to a declared
provider, capability, and action, making parts of this linkage visible to the
user.

Consequently, we evaluate our system from two perspectives: a technical benchmark
measuring contract coverage, non-regression, runtime enforcement, fault
localization, and overhead cost; alongside a user study examining whether
participants correctly identify the target object, observe handoffs, recognize
the responding agent, and understand the system's delegation of responsibilities
based on contract-driven visual cues. This empirical evaluation focuses
specifically on the implemented interaction, measuring task clarity,
predictability, guidance, and usability, rather than long-term trust or general
cross-domain applicability.

We address four central research questions:

\begin{description}
\item[RQ1---Coverage and non-regression.] Can multi-agent 3D environments
be adapted to end-to-end target--provider--capability--action chains while
maintaining existing ownership rules and routing logic under a fixed policy?

\item[RQ2---Fail-closed enforcement.] Does the implemented communication path
accept only valid local and single-hop requests, retain target and provider
evidence, and cleanly reject stale, ambiguous, unreachable, or conflicting
requests?

\item[RQ3---Added assurance and cost.] What cross-layer inconsistencies
can the contract catch and isolate beyond direct artifacts, and what are the
associated representation and performance overheads?

\item[RQ4---User comprehension.] How accurately do users perceive the chosen
target, the object--answer mapping, handoffs, and agent responsibilities, and
how predictable and supportive do they find the interaction?
\end{description}

This work provides four main contributions:

\begin{enumerate}
\item an interaction-centered failure taxonomy and executable contract binding
spatial selections to a single responsible provider, capability, action, and
evidence trace under a pinned model version;

\item a fail-closed system architecture spanning a Unity 3D client and
multi-agent backend featuring authoritative server-side resolution, staged
validation, strict one-hop handoff rules, and transactional response
verification;

\item a reproducible, API-free technical evaluation across three 3D
environments, including chain audits, non-regression comparisons with direct
wiring, routing probes, fault injection, and structural cost analysis; and

\item a user study evaluating target clarity, answer--target alignment,
handoff perception, agent role recognition, predictability, and user guidance.
\end{enumerate}

Our technical evaluation indicates that the contract maintains baseline system
behavior, strictly enforces one-hop policies, and detects cross-layer
discrepancies that native artifacts miss. The empirical user study complements
these findings: while users readily grasped object selection, target--answer
grounding, and handoff events, agent selection rationales and functional
boundaries were less intuitive. This demonstrates that while a contract can
reliably enforce backend responsibilities, additional interface cues are
necessary for full user clarity.

The complete prototype, model artifacts, evaluation scripts, raw data, study
materials, and reproduction guidelines are made available in the anonymous
repository. We do not claim that formal metamodels alone grant intelligence
to interfaces; rather, explicit relations constrain what systems interpret and
execute, reducing errors like wrong-target references or undeclared actions into
predictable, auditable bounds.

Section~\ref{sec:related} outlines related research on embodied agents, spatial
reference, conversational XR, model-driven engineering, and runtime
verification. Section~\ref{sec:contract} details the proposed contract and its
formal guarantees; Section~\ref{sec:implementation} describes its
implementation across Unity, the backend, and evidence tracing.
Sections~\ref{sec:evaluation} and~\ref{sec:results} present technical benchmarks
and user study findings. Section~\ref{sec:discussion} explores system trade-offs,
costs, and communication boundaries, while Sections~\ref{sec:limitations}
and~\ref{sec:conclusion} outline limitations and conclusions. Supplementary
materials are included in the appendix.
